% profiles/pennsylvania-state-university-main-campus.tex
% — Pennsylvania State University-Main Campus
% ============================================================================
% source: https://gradschool.psu.edu/academics/theses-and-dissertations
%         (the Fox Graduate School's Theses and Dissertations page)
% source: https://gradschool.psu.edu/assets/uploads/documents/Thesis-and-Dissertation-Handbook08272026.pdf
%         ("THESIS AND DISSERTATION HANDBOOK — Requirements and Guidelines for
%          the Preparation of Master's Theses and Doctoral Dissertations",
%          Office of Theses and Dissertations, J. Jeffrey and Ann Marie Fox
%          Graduate School. The filename carries its date: 2026-08-27, eight
%          days before this profile was written. 37 pp.)
% source: https://gradschool.psu.edu/academics/theses-and-dissertations/thesis-and-dissertation-templates
%         (the Graduate School's templates page — read for the accessibility
%          language quoted below, and for "Copyright ... This field is optional")
% accessed: 2026-09-05
% checked:  texlib-thesis-profile-round
%
% VERIFIED and set here: geometry, spacing, the title page, and the committee
% page. NOT set: chapter-opening margin (no such rule). No accessibility
% check-flag is raised — see the accessibility note, which is the Texas A&M
% pattern, not the Michigan one.
%
% READ THE MARGIN FIGURE FROM THE RENDERED PAGE, NOT THE EXTRACTED TEXT.
% `pdftotext` renders the handbook's vulgar fractions as U+FFFD, so p. 9 comes
% out as "A ?-inch or 1-inch margin", which reads naturally — and wrongly — as
% one-half. The page was rendered at 150 dpi and read as an image: the figure is
% THREE-QUARTERS. Anyone re-checking this profile must do the same.

\thesisinstitution{The Pennsylvania State University}

% --- Filing figures ----------------------------------------------------------
% "A 3/4-inch or 1-inch margin on all sides is acceptable, but keep in mind that
% a wider left margin (e.g., 1 1/2-inch) may be more appropriate for binding
% purposes."
%
% Two figures are acceptable and 1in is the larger, the class default, and the
% one an electronically filed eTD wants; 0.75in on all sides is equally in spec.
% The 1.5in left margin is explicitly a suggestion for BINDING ("may be more
% appropriate"), not a requirement, and is deliberately not encoded — Penn State
% files electronically, so a filed PDF should not carry a binding allowance.
\thesissetgeometry{left=1in, right=1in, top=1in, bottom=1in}

% "The text of an eTD may be single-, double- or one-and-a-half-spaced." All
% three are permitted and none is required, so this is a default, not a rule.
\thesissetspacing{double}

% No chapter-opening rule is published. Penn State does require "Begin each
% chapter on a new page", which the class already does, but states no deeper top
% margin there, so \thesissetchapteropening stays unset.

% --- Accessibility -----------------------------------------------------------
% NO CHECK-FLAG IS RAISED, deliberately. Penn State has clearly thought about
% this — but what it publishes is a statement about its own Word templates, in
% the indicative, with no modal verb anywhere: no "must", no "required", no
% named standard, no checker report to file. That is encouragement, not a
% filing rule, and \thesisrequiretagging is for rules. The handbook itself (37
% pp.) contains no document-accessibility requirement at all; its only
% "accessible" strings are about eTD access ("becomes accessible worldwide").
%
% The language is recorded verbatim so a later pass can see it was read and
% weighed, and can notice if the modal verb ever changes.
\thesisaccessibilityrequirement{Templates, not a mandate}
	{"Our templates were created with accessibility in mind. Predefined Styles
	 are applied to all headings, captions, and body text to ensure your
	 document meets accessibility standards." "Applying these Styles
	 consistently keeps your document accessible and professionally formatted."
	 Penn State names no standard, requires no tagged PDF and asks for no
	 checker report, so nothing here is enforced by this class.}

% --- Copyright line on the title page ------------------------------------------
% Penn State puts the copyright notice ON the title page, not on a page of its
% own, and the class has no macro for the year. \thesiscopyrightyear supplies it
% and defaults to empty, in which case the line is omitted entirely — which is
% correct, because the notice is optional twice over: "Copyright ... This field
% is optional" (templates page) and "You own the copyright on your work whether
% you choose to include this notice or not" (handbook).
%
% Set it in the document preamble, after \documentclass:  \thesiscopyrightyear{2027}
% The handbook's rule for the line: "use the word 'Copyright' or the symbol (c)
% (not both), followed by the year and your name." The symbol form is used, as
% in the handbook's own example page ("(c) 2023 John J. Doe").
\newcommand{\thesis@copyrightyear}{}
\newcommand{\thesiscopyrightyear}[1]{\renewcommand{\thesis@copyrightyear}{#1}}

% --- Title page --------------------------------------------------------------
% "The title page must appear exactly as the example title page of this guide."
% Reproduced from that example (handbook p. 30/31, rendered and read as an
% image), with the prose rules applied on top of it:
%
%   * "Type the title of the thesis/dissertation using capital letters
%     throughout. If it occupies more than one-line, double space between
%     lines."  -> \MakeUppercase, and the title block is double-spaced.
%   * "Master's candidates should use 'Thesis,' and doctoral candidates should
%     use 'Dissertation' on the title page."  -> \thesis@DocNoun, from \doctype.
%   * "On the date line, indicate the month and year of degree conferral, not
%     the date of the defense or the date you submit."  -> \graddate must be the
%     CONFERRAL date. Degrees are conferred only in May, August and December.
%   * "The title page does not show a page number, although it is actually page
%     i."  -> \thispagestyle{empty} with the no-footer geometry.
%
% The title is set bold because the example page sets it bold; the handbook only
% PERMITS bold there ("Boldface type may also be used on the title page"), so a
% document that wants it plain is not out of spec.
%
% AN ERROR IN THE HANDBOOK, left for the reviewer rather than papered over: the
% page headed "Doctoral Title Page Example" shows a MASTER'S title page — "A
% Thesis in ... Master of Science" — identical in substance to the page headed
% "Master's Title Page Example" on the next leaf. There is therefore no doctoral
% specimen. The doctoral wording below follows the handbook's PROSE rule quoted
% above, not the mislabelled example.
\renewcommand{\thesistitlepage}{%
	\loadgeometry{nopagenumber}%
	\begin{titlepage}%
		\thispagestyle{empty}\singlespacing\centering
		The Pennsylvania State University\par
		The J. Jeffrey and Ann Marie Fox Graduate School\par
		\vspace{4\baselineskip}
		% Double-spaced, per "If it occupies more than one-line, double space
		% between lines" -- so a two-line title breaks correctly by itself.
		\begingroup\doublespacing
		\thesis@tagtitle{Title}{{\bfseries\MakeUppercase{\@title}}}\par
		\endgroup
		\vspace{3\baselineskip}
		A \thesis@DocNoun{} in\par
		\vspace{\baselineskip}
		\thesis@program\par
		\vspace{2\baselineskip}
		by\par
		\vspace{\baselineskip}
		\@author\par
		\ifx\thesis@copyrightyear\@empty\else
			\vspace{2\baselineskip}
			\copyright\ \thesis@copyrightyear\ \@author\par
		\fi
		\vspace{3\baselineskip}
		Submitted in Partial Fulfillment\\
		of the Requirements\\
		for the Degree of\par
		\vspace{2\baselineskip}
		\thesis@degree\par
		\vspace{2\baselineskip}
		\thesis@date\par
		\vspace*{\fill}
	\end{titlepage}%
	\loadgeometry{normal}%
}

% --- Committee page ------------------------------------------------------------
% Penn State calls it the COMMITTEE PAGE, not an approval page, and it is part of
% the filed document: "The committee page is page ii and appears immediately
% before the Abstract." "The committee page must appear exactly as example
% committee page of this guide."
%
% NO SIGNATURE LINES, AND NO PHYSICAL SIGNATURES. The page carries typed names
% and titles only; approval is collected electronically through the eTD system
% ("All electronic signatories on a doctoral dissertation must be members of
% Penn State's Graduate Faculty..."). Drawing a rule here would be wrong.
%
% Opening sentence, from the example: "The dissertation of John J. Doe was
% reviewed and approved by the following:" — so it takes the document noun,
% lowercase, and the author's name.
%
% Per-member shape, from the example plus the prose: the name, then the
% professorial title, then any role(s), each on its own line, left-aligned.
% "Be sure that all committee members are identified by their correct
% professorial titles. ... Do not use such designations as 'Ph.D.' or 'Dr.' on
% the committee page." So a document filing here writes, e.g.,
%   \committeemember{Charles L. Merkle}{Professor of Mechanical Engineering}
% and puts the ROLE ("Dissertation Advisor", "Chair of Committee") on the same
% second argument, separated by \\ if there are two — the class's committee list
% carries exactly two fields.
%
% NOTE ON THE EXAMPLE'S OWN TYPOGRAPHY: in the handbook's rendered example the
% advisor's first entry reads "Professor of Mechanical Engineering Dissertation
% Advisor" on ONE line, with "Chair of Committee" beneath. The prose says roles
% go "under the professorial title", so that run-on is treated as a slip in the
% specimen and the roles are set on their own lines here.
%
% THE PAGE IS NUMBERED. Unlike the title page, the committee page shows "ii", so
% this page keeps the normal geometry and page style rather than suppressing the
% number. For it to come out as a lowercase roman numeral the document must call
% \thesisapprovalpage inside the front matter -- i.e. after \frontmatter -- which
% is the document's business, not the profile's.
\renewcommand{\thesisapprovalpage}{%
	\clearpage
	\loadgeometry{normal}\singlespacing
	\thesis@tagtitle{H1}{{\color{white}\tiny Committee Page}}\par
	\begingroup
	\raggedright
	The \thesis@docnoun{} of \@author{} was reviewed and approved by the
	following:\par
	\renewcommand{\thesis@memberline}[2]{%
		\vskip 1.5em
		##1\par
		\ifx\relax##2\relax\else ##2\par\fi
	}%
	\thesis@committee
	\endgroup
	\doublespacing\clearpage
}
